\section{Introduction}\label{sec:introduction}

Let $p$ be a prime, given in binary, and let $f\in\Fp[x]$ be a
nonzero polynomial, given by its dense coefficient list. Write
$n=\deg f$ and $L=\lceil\log_2 p\rceil$. Complete factorization means
computing the leading coefficient $c$ and distinct monic irreducible
polynomials $g_i$, with positive integer multiplicities $e_i$, such that
\[
 f=c\prod_i g_i^{e_i}.
\]
The empty product is permitted when $n=0$. The complexity parameter is
$(n+1)L$, so an algorithm polynomial in $p$ does not give a polynomial-time
algorithm in the prescribed representation.

Finite-field factorization is a basic constructive problem in algebra.
It also exposes a persistent distinction between randomized and
deterministic computation: a random element of a product of finite fields
often separates its components, whereas constructing a separator uniformly
can require additional arithmetic information. The present paper isolates
one form of that information and gives a deterministic splitting procedure
once it is available. To supply it uniformly, we use the zero-free theorem
for finite-order Hecke $L$-functions in the companion paper
\cite[Theorem~1.2]{OpenAIPrimitive26}. We state that input precisely in
Theorem~\ref{analytic:hecke} and prove its auxiliary-prime consequence
here. The analytic theorem itself belongs to the companion.

\begin{theorem}\label{thm:factorization}
There is a uniform deterministic algorithm that, given a prime $p$ in
binary and a nonzero dense polynomial $f\in\Fp[x]$ of degree $n$, outputs
its leading coefficient and its distinct monic irreducible factors with
their multiplicities in
\[
 O\bigl(((n+1)\lceil\log_2p\rceil)^{10^{12}}\bigr)
\]
bit operations, with absolute implied constant. A constant polynomial
returns an empty factor list. The algorithm uses no randomness,
integer-factorization oracle, primitive-root oracle, or GRH assumption.
\end{theorem}

The exponent is deliberately generous. The point is one fixed polynomial
bound for all degrees and prime characteristics. We separate the
algebraic construction from its analytic input in the following reduction.

\subsection{The auxiliary-prime reduction}

Put
\begin{equation}\label{eq:B}
 B=20+(n+1)(L+1).
\end{equation}
For a prime $q\le n$, call a rational prime $\ell$ an
\emph{auxiliary prime for $(p,q)$} if
\begin{equation}\label{eq:auxiliary}
 \ell\notin\{2,3,p,q\},\qquad \ell\equiv1\pmod{12q},\qquad
 p^{(\ell-1)/q}\not\equiv1\pmod\ell.
\end{equation}
Thus $p$ is not a $q$th power modulo $\ell$. The congruence modulo $12q$
is convenient for the analytic construction in Section~\ref{sec:analytic};
the algebraic construction only requires the corresponding congruence
modulo $q$.

\begin{theorem}\label{thm:reduction}
There is a uniform deterministic algorithm with the following guarantee.
Its input is a prime $p$ and a nonzero dense polynomial $f\in\Fp[x]$.
If $p>B^{200000}$ and $n\ge2$, its input also includes one auxiliary
prime $\ell_q$ satisfying \eqref{eq:auxiliary} for each prime $q\le n$.
It outputs the complete factorization of $f$, including multiplicities.
Let $E=\max(2,\max_q\ell_q)$, with $E=2$ when no auxiliary primes are
required. For an absolute constant $C$, the algorithm uses
$O((B+E)^C)$ bit operations. It uses neither randomness nor an integer
factorization or primitive-root oracle.
\end{theorem}

The dependence on $E$ is on the numerical value of the auxiliary primes,
not merely their bit lengths. Consequently Theorem~\ref{thm:reduction}
becomes a polynomial-time factorization theorem when these primes are
bounded by a fixed power of $B$. Upward search, trial division, and modular
exponentiation then construct them within the same kind of bound. We prove
this consequence from the cited zero-free theorem in
Section~\ref{sec:analytic}. The reduction in
Theorem~\ref{thm:reduction} does not use that theorem.

\subsection{History and the role of the construction}
\label{sec:history}

Berlekamp's algorithms put linear algebra at the center of finite-field
factorization: Frobenius-fixed elements in the quotient algebra encode its
irreducible factors, and polynomial greatest common divisors turn suitable
such elements into factors~\cite{Berlekamp67,Berlekamp70}.  Randomized
algorithms, including Cantor--Zassenhaus~\cite{CantorZassenhaus81}, obtain
expected polynomial running time by finding separating elements with random
choices.  The difficulty studied here is to make every such choice
deterministically, with a bound polynomial in both the degree and the bit
length of the characteristic.

Several deterministic results delineate this difficulty.  Shoup's algorithm
has polynomial dependence on the degree and a
$p^{1/2+o(1)}$ dependence on the characteristic~\cite{Shoup90}.
Under the generalized Riemann hypothesis, R\'onyai obtained polynomial
bounds when the number of irreducible factors is bounded~\cite{Ronyai88},
and Evdokimov obtained a bound
$(n^{\log n}\log p)^{O(1)}$ for general degree~\cite{Evdokimov94}.
The latter is quasipolynomial in $n$. The reduction proved here, combined
with the fixed-width zero-free strip that follows immediately from Hecke
GRH, already gives a polynomial bound in both parameters under GRH.
For Theorem~\ref{thm:factorization}, the companion supplies that strip
without GRH. In particular, the hypotheses and
parameters of a deterministic result matter: removing random choices
alone does not establish a uniform polynomial bound in $(n+1)L$.

Geometric methods already have an important place in this history.
Schoof's elliptic-curve method computes a square root of a fixed integer
$a$ modulo a varying prime, when one exists, in time polynomial in $\log p$; the uniform
bound in the two parameters is polynomial in $|a|$ and $\log p$, rather
than in their binary lengths~\cite[Section~4]{Schoof85}.
Pila's computation of Frobenius on abelian varieties yields deterministic
construction of roots of unity of fixed odd prime order modulo primes
congruent to $1$ modulo that order, with a running-time exponent that may
depend on the order and on the geometric
presentation~\cite[Theorems~A and~D]{Pila90}.  These results motivate
careful accounting for the varying curve and prime in the present
construction.  We must bound the degrees of field extensions, the sizes
of divisor representations, and the dimensions of all linear systems by
fixed powers of the original input parameters.

For a single prime not dividing the discriminant of an integral lift,
R\'onyai obtained, under GRH, deterministic bounds involving the degree
of that lift's splitting field~\cite{Ronyai92}.
A different kind of uniformity arises when reductions modulo many primes
are processed together.  Altman's unconditional algorithm takes
a monic irreducible polynomial in $\mathbb Z[x]$ of coefficient height
at most $H\ge2$, with splitting-field degree $m$, and factors its
reductions modulo every prime below $X$ in total time
\[
 \pi(X)(m\log H)^{O(1)}\log^5X.
\]
This is an amortized result over primes~\cite[Theorem~1.1]{Altman2025}.
The parameter $m$ can be as large as the factorial of the polynomial's
degree.  Thus this result does not give the worst-case single-input
bound sought here.  Altman's introduction also provides a recent account
of the general deterministic problem and its remaining obstacles.

The algebraic techniques used below have substantial precedents.
Computation at unknown roots by working in a quotient algebra and
splitting it when a zero test distinguishes components is an instance of
dynamic evaluation~\cite{DellaDoraDicrescenzoDuval85,Duval94}.
The digit calculations in primary multiplicative groups are the
prime-power part of the Pohlig--Hellman algorithm~\cite{PohligHellman78}.
R\'onyai already proved that a completely split polynomial of even degree
can be split in polynomial time when a quadratic nonresidue is
supplied~\cite{Ronyai88}; we include an explicit tournament proof suited
to our representation.  The odd-degree construction uses cyclic curves
$Y^q=F(X)$.  The cyclotomic action on their Jacobians and descriptions of
its first torsion layer by ramification divisors are familiar from
explicit descent, notably the work of Poonen and
Schaefer~\cite[Sections~4 and~6]{PoonenSchaefer97}.

Our effective treatment of divisors also builds on established
computational geometry.  Ideal representations on superelliptic curves,
Riemann--Roch algorithms, and finite-dimensional linear algebra for
Jacobian arithmetic appear in
\cite{GPS02,Hess02,KhuriMakdisi04}.  For norm equations, we use the
relation between cyclic algebras and norms, and the approach through
maximal orders and minimal left ideals developed in explicit
matrix-algebra computations~\cite{GilleSzamuely06,IRS12,IKR18}.
The general function-field algorithm of Ivanyos, Kutas, and R\'onyai
uses a finite-field polynomial-factorization oracle.  Here the cyclic
presentation gives explicit local orders, and evaluation of their global
sections at a rational point gives a rank-one idempotent by linear
algebra.  The structural explanation uses the splitting of vector
bundles on the projective line~\cite{Grothendieck57,HazewinkelMartin82}.

The substantive step of the present reduction is the odd-degree
separator, together with a uniform implementation of the divisor-class
divisions it requires.  A filtration of lifted ramification classes
keeps the necessary field extensions of polynomial degree.  A lattice
of divisors at infinity then forces two ramification points to receive
different labels.  We prove this mechanism and its algorithms directly;
the earlier results just discussed provide their mathematical context.
The separate step of finding sufficiently small auxiliary primes is
isolated in Section~\ref{sec:analytic}, which is the only place the
companion's analytic theorem enters.

\subsection{Proof strategy}

The Berlekamp algebra reduces complete factorization to splitting a monic
square-free polynomial $F$ all of whose roots already belong to $\Fp$.
Write $N=\deg F\ge2$. The unknown roots are represented simultaneously
by the finite algebra $\Fp[T]/(F)$; an operation whose outcomes differ
between components immediately produces a factor by a greatest common
divisor. We can therefore describe computations at an unspecified root
and carry them out without first finding that root.

The auxiliary primes construct the multiplicative information needed in
these computations. For each odd prime $q$, they give a field
$K_q=\Fp(\zeta_q)$ and a generator of the full $q$-primary subgroup of
$K_q^*$. The entry for $q=2$ is a generator of the full $2$-primary
subgroup of $\Fp^*$. Here the $q$-primary subgroup consists of the
elements whose orders are powers of $q$. These entries permit extraction
of promised $q$th roots in finite algebras by small-prime discrete
logarithms. Their construction proceeds in increasing $q$: the only
factorization needed to make the entry for $q$ has degree $q-1$.

When $N$ is even, differences between pairs of roots and the $2$-primary
entry orient the complete graph on the roots. Opposite differences have
opposite orientations. The vertex scores cannot all be equal when $N$
is even, so they provide a separator. The more substantial construction
handles odd $N$.

Choose an odd prime $q\mid N$, and write $r=v_q(N)$. Over an explicitly
constructed field $k$ of degree at most $N(q-1)$ over $\Fp$, consider
the smooth projective curve with affine equation
\[
 C:\quad Y^q=F(X).
\]
Its ramification points are $P_i=(x_i,0)$, where the $x_i$ are the
unknown roots of $F$. There are $q$ rational points at infinity.
The automorphism $\sigma(X,Y)=(X,\zeta_qY)$ acts on degree-zero divisor
classes; write $\lambda=1-\sigma$. Degree-zero divisors supported at
infinity form a lattice on which $\lambda$ is injective. In contrast,
$\lambda$ is surjective on the geometric Jacobian.

This difference is the source of the splitting argument. We lift each
class $[P_i-\infty_0]$ repeatedly through $\lambda$, keeping track of
an infinity divisor at every stage. A filtration argument shows that all
the classes required for these lifts are rational over the field $k$;
the extension degree stays polynomial even though the number of possible
classes is large. Summing the lifted classes and applying explicit label
tests forces a distinction between ramification points. If every test
gave one common label, a fixed nonzero infinity divisor would become
divisible by one more power of $\lambda$ than its lattice valuation
allows.

The geometric existence argument is accompanied by a constructive
division algorithm. A promised $\lambda$-division reduces to a norm
equation in $k(C)/k(X)$. We solve that equation by constructing maximal
orders in a split cyclic algebra and computing their global sections
on the projective line. A rank-one idempotent in this algebra yields
the norm solution. All computations use polynomial and matrix arithmetic;
Riemann--Roch reduction controls divisor sizes, and arithmetic circuits
store principal functions arising from large infinity multiplicities.
Figure~\ref{fig:route} records the roles of these ingredients.

\begin{figure}[t]
\centering
\begin{tikzpicture}[font=\small,>=Latex,
 box/.style={draw,rounded corners=2pt,align=center,text width=5.2cm,
 minimum height=11mm,inner sep=5pt},node distance=9mm and 12mm]
 \node[box,dashed] (analytic) {Uniform Hecke zero-free theorem\\from the companion};
 \node[box,below=of analytic] (primes) {Polynomially bounded auxiliary primes};
 \node[box,below=of primes] (table) {Finite fields and\\primary generators};
 \node[box,below left=12mm and -12mm of table] (norm)
   {Norm equations and\\divisor-class division};
 \node[box,below right=12mm and -12mm of table] (geometry)
   {Rational lifts and the\\infinity-lattice obstruction};
 \node[box,below=28mm of table] (split) {A factor of a\\totally split polynomial};
 \node[box,below=of split] (factor) {Complete factorization\\with multiplicities};
 \draw[->,dashed] (analytic)--(primes);
 \draw[->] (primes)--(table);
 \draw[->] (table)--(norm);
 \draw[->] (table)--(geometry);
 \draw[->] (norm)--(split);
 \draw[->] (geometry)--(split);
 \draw[->] (split)--(factor);
\end{tikzpicture}
\caption{The proof separates auxiliary-prime existence from the algebraic
reduction. The dashed implication uses the cited analytic theorem;
the remaining constructions are proved in the paper. The two middle
branches are the effective and geometric parts of the odd-degree splitter.}
\label{fig:route}
\end{figure}

The technical results are uniform in the degree of the curve and the
prime $q$. This uniformity is essential: a running-time exponent depending
on the genus would not suffice for Theorem~\ref{thm:reduction}.
We give the finite-field tools first, then the auxiliary table, the
geometry, and the effective divisor and norm calculations. The splitter
and the factorization driver assemble these ingredients. The final two
sections prove the bit bounds and the analytic implication.
